\part{散度函数与对偶平坦黎曼结构}\label{part:amari-partI}
\addcontentsline{toc}{part}{Part I: 散度函数与对偶平坦黎曼结构}

本部分对应 Amari《Information Geometry and Its Applications》Part I (Chapter 1-4)，介绍信息几何的基础理论。

\section{Part I 概述}

Part I 建立了信息几何的核心数学框架，包括：

\begin{enumerate}
    \item \textbf{Chapter 1}: 流形、散度与对偶平坦结构
    \item \textbf{Chapter 2}: 指数族与混合族
    \item \textbf{Chapter 3}: 不变几何与 Fisher 信息
    \item \textbf{Chapter 4}: $\alpha$-几何与 Tsallis $q$-熵
\end{enumerate}

\section{核心概念}

本Part的核心概念包括：

\begin{itemize}
    \item \textbf{散度（Divergence）}：概率分布之间的非对称"距离"
    \item \textbf{Bregman 散度}：由凸函数诱导的散度
    \item \textbf{指数族}：$p(x|\theta) = \exp(\theta \cdot u(x) - \psi(\theta))$
    \item \textbf{Fisher 信息}：概率分布空间的不变度量
    \item \textbf{对偶平坦结构}：由 Legendre 变换联系的两组仿射坐标
\end{itemize}

\noteinfo{
    Part I 的核心洞见是：\textbf{散度函数诱导了概率分布流形上的几何结构}，这为后续的统计推断提供了统一的几何框架。
}
